\section{Formal Grammar Elaboration Rules}
\label{sec:appendix-grammar}

The TKF family of evolutionary models---TKF91, TKF92, MixDom,
the TKF Structure Tree, and the TKF Genome---describes the joint
evolution of biological sequences subject to insertions, deletions,
and substitutions.  Despite their apparent diversity, all these models
share a common constructive pattern: they begin with a simple weighted
context-free grammar (WCFG) for a geometrically distributed number of
links, and then systematically elaborate that grammar through a
series of formal transformations.

In the TKF91 model~\cite{ThorneEtAl91}, each link carries a single character evolving by a
continuous-time Markov chain (CTMC).  TKF92~\cite{ThorneEtAl92}
extends this by replacing each character with a geometrically
distributed fragment of characters.  The MixDom model nests a TKF92
process inside a TKF91 process, decorating links with mixtures of
domain types, fragment types, and substitution classes.  The TKF
Structure Tree~\cite{Holmes2004} uses stochastic context-free grammar
(SCFG) recursion to model RNA secondary structure with stems (emitting
paired characters left and right) and loops.  The TKF Genome extends
these ideas to entire genomes with coding sequences, introns, RNA structures,
and conserved elements.

Making the elaboration steps explicit and composable has several benefits:
\begin{enumerate}[label=(\roman*)]
\item \textbf{Correctness}: each transformation can be verified independently,
  rather than checking a large monolithic grammar.
\item \textbf{Modularity}: new models (e.g., RNA models with basepair stacking,
  codon models with reading-frame-aware indels) can be constructed by
  composing well-understood building blocks.
\item \textbf{Automation}: the transformations are sufficiently formal
  to be implemented as software operations on grammar objects,
  enabling automatic derivation of dynamic programming algorithms
  from high-level model specifications.
\end{enumerate}

This appendix defines seven elaboration rules and the associated
null-state management procedures, then shows how each known
TKF-family model arises as a specific sequence of elaborations.

\subsection{Base Grammar}

\subsubsection{Weighted Context-Free Grammars}

\begin{definition}[Weighted Context-Free Grammar]
A \emph{weighted context-free grammar} (WCFG) is a tuple
$\WCFG = (\NT, \TM, \PR, S, \WF)$ where:
\begin{itemize}
\item $\NT$ is a finite set of nonterminal symbols.
\item $\TM$ is a finite set of terminal symbols, disjoint from $\NT$.
\item $S \in \NT$ is the start symbol.
\item $\PR$ is a finite set of production rules of the form $X \to \alpha$
  where $X \in \NT$ and $\alpha \in (\NT \cup \TM)^*$.
\item $\WF : \PR \to \mathbb{R}_{\geq 0}$ assigns a nonneg weight to each production.
\end{itemize}
The grammar is \emph{proper} if for every nonterminal $X$,
the weights of all productions with left-hand side $X$ sum to 1:
$\sum_{(X \to \alpha) \in \PR} \WF(X \to \alpha) = 1$.
In a proper WCFG, weights are probabilities and the grammar defines
a stochastic context-free grammar (SCFG).
\end{definition}

\begin{definition}[Elaboration Rule]
An \emph{elaboration rule} (or grammar transformation) is a map
$\mathcal{E} : \WCFG \to \WCFG'$ that takes a WCFG
and a set of elaboration parameters, and produces a new WCFG.
An elaboration is \emph{validity-preserving} if it maps
proper grammars to proper grammars.
\end{definition}

We distinguish between the \emph{single-sequence grammar}
(describing the stationary distribution over sequences)
and the \emph{pair grammar} (describing the joint distribution
over ancestor--descendant sequence pairs).  Most elaborations
operate on the single-sequence grammar; the Evolution elaboration
(Section~\ref{sec:evolution-appendix}) converts a single-sequence grammar
into a pair grammar.

\subsubsection{The Link Grammar}
\label{sec:base-appendix}

The fundamental building block of all TKF-family models is a grammar
generating a geometrically distributed number of ``links.''
This grammar arises from the stationary distribution of the linear
birth-death-immigration (BDI) process with per-capita birth rate
$\insrate$, per-capita death rate $\delrate > \insrate$, and
immigration rate $\immrate = \insrate$.

\begin{definition}[Link Grammar]
\label{def:link-grammar-appendix}
The \emph{link grammar} $\WCFG_{\mathrm{link}}(\kappa)$ with
parameter $\kappa = \insrate/\delrate \in [0,1)$ is the WCFG
with nonterminals $\{\texttt{IMM}, \texttt{MOR}\}$, start symbol
$\texttt{IMM}$, and productions:
\begin{align}
\texttt{IMM} &\to \texttt{MOR}\;\texttt{IMM} && \text{weight } \kappa \label{eq:imm-extend} \\
\texttt{IMM} &\to \epsilon && \text{weight } 1-\kappa \label{eq:imm-end} \\
\texttt{MOR} &\to \texttt{MOR}\;\texttt{MOR} && \text{weight } \kappa \label{eq:mor-extend} \\
\texttt{MOR} &\to \epsilon && \text{weight } 1-\kappa \label{eq:mor-end}
\end{align}
\end{definition}

\begin{remark}
In this grammar, $\texttt{IMM}$ (the ``immortal link'') generates
a sequence of $n \sim \geomdist(\kappa)$ mortal links.
Each $\texttt{MOR}$ can recursively generate further mortal links;
the recursive self-loop in~\eqref{eq:mor-extend} reflects the
offspring-generating property of the BDI process.
The distinction between $\texttt{IMM}$ and $\texttt{MOR}$
captures the different roles: the immortal link corresponds to the
BDI regime $\immrate = \insrate, X(0) = 0$ (immigration from nothing),
while mortal links correspond to $\immrate = 0, X(0) = 1$
(a single founder that can die).
\end{remark}

\begin{remark}
The grammar in Definition~\ref{def:link-grammar-appendix} generates only
the empty string $\epsilon$, since $\texttt{MOR}$ has no terminal-producing
rules.  The elaboration rules below will add terminal emissions
(characters, character pairs, etc.) to the mortal links.
\end{remark}

\begin{proposition}
The link grammar $\WCFG_{\mathrm{link}}(\kappa)$ is proper for
any $\kappa \in [0,1)$.  Under the start symbol $\texttt{IMM}$,
the number of $\texttt{MOR}$ expansions
before reaching $\epsilon$ is distributed as $\geomdist(\kappa)$
(with support $\{0,1,2,\ldots\}$ and mean $\kappa/(1-\kappa)$).
\end{proposition}

\begin{proof}
For $\texttt{IMM}$: the productions~\eqref{eq:imm-extend}
and~\eqref{eq:imm-end} have weights $\kappa$ and $1-\kappa$,
summing to 1.  Similarly for $\texttt{MOR}$.
The number of $\texttt{MOR}$ nonterminals generated by
$\texttt{IMM}$ before choosing $\epsilon$
is geometric with parameter $\kappa$ by the standard
geometric series argument.
\end{proof}

\subsection{Elaboration Rules}

We now define each elaboration rule as a formal grammar transformation.
For each rule, we specify:
\begin{itemize}
\item The \emph{input} grammar fragment (which productions are targeted).
\item The \emph{output} grammar fragment (the replacement productions).
\item The \emph{parameters} introduced by the elaboration.
\item The \emph{validity conditions} under which the transformation preserves properness.
\end{itemize}

\subsubsection{CTMC Expansion}
\label{sec:ctmc-expansion-appendix}

The most basic elaboration decorates each link with a character
(or tuple of characters) that evolves according to a finite-state
CTMC.  This is the step that takes the bare link grammar to a model
with observable sequences.

\begin{definition}[Emission Type]
An \emph{emission type} specifies where terminal symbols appear
relative to the recursive expansion of a nonterminal:
\begin{itemize}
\item \textbf{Left-emission}: terminal appears to the left of the recursive part.
  Production form: $X \to c\;\alpha$ where $c \in \TM$.
\item \textbf{Right-emission}: terminal appears to the right.
  Production form: $X \to \alpha\;c$.
\item \textbf{LR-emission}: terminals appear on both sides.
  Production form: $X \to c_L\;\alpha\;c_R$ where $c_L, c_R \in \TM$.
  This is the form needed for RNA basepair models (stems).
\end{itemize}
\end{definition}

\begin{definition}[CTMC Expansion]
\label{def:ctmc-appendix}
Let $\WCFG$ be a WCFG containing a nonterminal $X$ with an
$\epsilon$-generating production used for link termination.
Let $\alphabet$ be a finite alphabet, $\eqm$ a probability distribution
over $\alphabet$, and $\exch$ a rate matrix with stationary distribution $\eqm$.
The \emph{left CTMC expansion} of $X$ with parameters $(\alphabet, \eqm, \exch)$
replaces every production $X \to \alpha$ (where $\alpha \neq \epsilon$)
with the family of productions:
\begin{equation}
X \to c\;\alpha \qquad \text{weight } \WF(X \to \alpha) \cdot \eqm_c
\quad \text{for each } c \in \alphabet
\end{equation}
The $\epsilon$-production $X \to \epsilon$ is left unchanged.

The \emph{right CTMC expansion} replaces $X \to \alpha$ ($\alpha \neq \epsilon$) with:
\begin{equation}
X \to \alpha\;c \qquad \text{weight } \WF(X \to \alpha) \cdot \eqm_c
\quad \text{for each } c \in \alphabet
\end{equation}

The \emph{LR CTMC expansion} with alphabet $\alphabet \times \alphabet$
and equilibrium distribution $\eqm(c_L, c_R)$ replaces
$X \to \alpha$ ($\alpha \neq \epsilon$) with:
\begin{equation}
X \to c_L\;\alpha\;c_R \qquad \text{weight } \WF(X \to \alpha) \cdot \eqm(c_L, c_R)
\quad \text{for each } (c_L, c_R) \in \alphabet \times \alphabet
\end{equation}
\end{definition}

\begin{proposition}
Left, right, and LR CTMC expansion are validity-preserving:
if $\WCFG$ is proper, then so is the elaborated grammar.
\end{proposition}

\begin{proof}
For left CTMC expansion: the total weight of productions with
LHS $X$ becomes
$\sum_{c \in \alphabet} \sum_{\alpha \neq \epsilon} \WF(X \to \alpha) \cdot \eqm_c
+ \WF(X \to \epsilon)
= \sum_{\alpha \neq \epsilon} \WF(X \to \alpha) \cdot 1 + \WF(X \to \epsilon)
= \sum_\alpha \WF(X \to \alpha) = 1$.
The right and LR cases are analogous.
\end{proof}

\begin{example}[TKF91]
\label{ex:ctmc-tkf91-appendix}
Applying left CTMC expansion to $\texttt{MOR}$ in the link grammar
$\WCFG_{\mathrm{link}}(\kappa)$ with alphabet $\alphabet$
and equilibrium distribution $\eqm$ yields:
\begin{align*}
\texttt{IMM} &\to \texttt{MOR}\;\texttt{IMM} && \kappa \\
\texttt{IMM} &\to \epsilon && 1-\kappa \\
\texttt{MOR} &\to c\;\texttt{MOR}\;\texttt{MOR} && \kappa \cdot \eqm_c \quad (c \in \alphabet)\\
\texttt{MOR} &\to c && (1-\kappa) \cdot \eqm_c \quad (c \in \alphabet)
\end{align*}
This is the stationary (single-sequence) grammar for TKF91:
it generates a geometric number of links, each carrying a character
drawn i.i.d.\ from $\eqm$.
\end{example}

\begin{example}[TKF Structure Tree stems]
\label{ex:ctmc-lr-appendix}
In the TKF Structure Tree, stems use LR CTMC expansion.
The nonterminal $S$ (``stem'') generates basepairs:
\begin{align*}
S &\to c_L\;S\;c_R && (1-\kappa_S) \cdot \eqm_S(c_L, c_R) \\
S &\to L && \kappa_S
\end{align*}
where $L$ is the loop nonterminal (which uses left CTMC expansion),
and $\eqm_S(c_L, c_R)$ is the joint equilibrium distribution over
basepairs.  The LR emission enables the grammar to generate
palindromic structures characteristic of RNA secondary structure.
\end{example}

\subsubsection{Fragment Expansion}
\label{sec:fragment-appendix}

Fragment expansion takes TKF91 to TKF92 by replacing each single-character
link with a geometrically distributed sequence of characters.

\begin{definition}[Fragment Expansion]
\label{def:fragment-appendix}
Let $\WCFG$ be a WCFG with a nonterminal $\texttt{MOR}$ representing
a mortal link.  The \emph{fragment expansion} with parameter
$\ext \in [0,1)$ (the extension probability) replaces $\texttt{MOR}$
with three nonterminals $\texttt{MOR\_S}$ (fragment start),
$\texttt{MOR\_X}$ (fragment extend), and $\texttt{MOR\_E}$
(fragment end), defined as follows.

Every production in $\WCFG$ that references $\texttt{MOR}$
on its right-hand side is updated to reference $\texttt{MOR\_S}$ instead.
Then:

\medskip
\noindent\textbf{Before} (link with single terminal slot):
\begin{align*}
\texttt{MOR} &\to [\text{terminal}]\;\alpha && \WF_{\text{orig}}
\end{align*}

\noindent\textbf{After} (link with geometric fragment):
\begin{align*}
\texttt{MOR\_S} &\to \texttt{MOR\_X} && 1 \\
\texttt{MOR\_X} &\to [\text{terminal}]\;\texttt{MOR\_X} && \ext \\
\texttt{MOR\_X} &\to [\text{terminal}]\;\texttt{MOR\_E} && 1 - \ext \\
\texttt{MOR\_E} &\to \alpha && \WF_{\text{orig}}
\end{align*}

Here $[\text{terminal}]$ denotes whatever terminal emission was
associated with $\texttt{MOR}$ (a single character from a CTMC expansion,
or a character pair, etc.), and $\alpha$ denotes the rest of the
original production's right-hand side (the recursive continuation).

More precisely, if $\texttt{MOR}$ had productions
$\texttt{MOR} \to c\;\texttt{MOR}\;\texttt{MOR}$ (weight $\kappa \cdot \eqm_c$)
and $\texttt{MOR} \to c$ (weight $(1-\kappa) \cdot \eqm_c$)
after CTMC expansion, the fragment expansion produces:
\begin{align*}
\texttt{MOR\_S} &\to \texttt{MOR\_X} && 1 \\
\texttt{MOR\_X} &\to c\;\texttt{MOR\_X} && \ext \cdot \eqm_c \quad (c \in \alphabet) \\
\texttt{MOR\_X} &\to c\;\texttt{MOR\_E} && (1-\ext) \cdot \eqm_c \quad (c \in \alphabet) \\
\texttt{MOR\_E} &\to \texttt{MOR\_S}\;\texttt{MOR\_S} && \kappa \\
\texttt{MOR\_E} &\to \epsilon && 1-\kappa
\end{align*}
where $\texttt{MOR\_E}$ takes over the inter-link continuation logic
from the original $\texttt{MOR}$.
\end{definition}

\begin{proposition}
Fragment expansion is validity-preserving.
The expected number of terminals per fragment is $1/(1-\ext)$.
\end{proposition}

\begin{remark}
Fragment expansion must be applied \emph{after} CTMC expansion
(or simultaneously), because it assumes the link already has
terminal emissions.  If applied to a bare link grammar (no terminals),
the result would have fragments of $\epsilon$'s, which is degenerate.
\end{remark}

\begin{example}[TKF92]
TKF92 = link grammar $\to$ (fragment expansion with parameter $\ext$)
$\to$ (CTMC expansion with $\alphabet, \eqm, \exch$).
Equivalently, CTMC expansion first, then fragment expansion.
Both orderings yield the same grammar, because fragment expansion
simply wraps the terminal emission in a geometric self-loop.
The resulting grammar generates sequences where each link produces
a fragment of $K \sim \geomdist(\ext)$ characters from $\eqm$.
\end{example}

\subsubsection{Mixture Expansion}
\label{sec:mixture-appendix}

Mixture expansion decorates a link with a latent categorical variable
whose value determines subsequent model parameters.

\begin{definition}[Mixture Expansion]
\label{def:mixture-appendix}
Let $\WCFG$ be a WCFG containing a nonterminal $X$.
Let $\{1, \ldots, K\}$ be a finite set of mixture components
with weights $p_1, \ldots, p_K > 0$ satisfying $\sum_k p_k = 1$.
The \emph{mixture expansion} of $X$ with components $K$
and weights $(p_k)$ replaces $X$ with $K$ new nonterminals
$X_1, \ldots, X_K$ and modifies all productions referencing $X$
as follows.

Every production $Y \to \alpha\;X\;\beta$ in $\WCFG$
(where $Y \neq X$ and $\alpha, \beta \in (\NT \cup \TM)^*$)
is replaced by $K$ productions:
\begin{equation}
Y \to \alpha\;X_k\;\beta \qquad \text{weight } \WF(Y \to \alpha X \beta) \cdot p_k
\quad \text{for } k = 1, \ldots, K
\end{equation}

The productions of $X$ itself are copied to each $X_k$:
for each production $X \to \gamma$, create
\begin{equation}
X_k \to \gamma_k \qquad \text{weight } \WF(X \to \gamma)
\end{equation}
where $\gamma_k$ is $\gamma$ with any self-references to $X$ replaced
by $X_k$ (maintaining the component assignment within a link).

Each $X_k$ may then undergo different subsequent elaborations
(e.g., different CTMC parameters, different fragment extension rates).
\end{definition}

\begin{proposition}
Mixture expansion is validity-preserving.
\end{proposition}

\begin{proof}
For $Y$: the total weight of productions referencing $X_k$ is
$\sum_k \WF(\cdot) \cdot p_k = \WF(\cdot) \cdot 1$.
For each $X_k$: the production weights are copied from $X$, hence sum to 1.
\end{proof}

\begin{remark}
In principle, the mixture component selector could itself evolve
via a CTMC across evolutionary time, provided that nested expansions
are resampled from equilibrium upon a change of component.
In this appendix, we restrict attention to the simpler case where
the component assignment is fixed throughout the lifetime of a link
(i.e., it is sampled once at birth and does not change).
\end{remark}

\begin{example}[MixDom: three levels of mixture]
In the MixDom model, mixture expansion is applied at three levels:
\begin{enumerate}
\item \textbf{Domain mixture}: each top-level link is assigned a
  domain type $\dom \sim \catdist(\domdist_1,\ldots,\domdist_\ndom)$, determining
  $(\insrate_\dom, \delrate_\dom)$ for the nested TKF92 process.
\item \textbf{Fragment process}: within domain $\dom$, the initial
  fragment type is $\frag \sim \catdist(\fragdist_{\dom 1},\ldots,\fragdist_{\dom\nfrag})$.
  Subsequent fragments are drawn
  from the $\nfrag \times \nfrag$ transition matrix $\ext^{(\dom)}_{\srcfrag\destfrag}$
  (the $\nfrag = 1$ case reduces to IID geometric extension).
\item \textbf{Site class mixture}: within fragment state $\frag$ of domain $\dom$,
  each character position is assigned a class $\class \sim \catdist(\classdist_{\dom\frag 1},\ldots,\classdist_{\dom\frag\nclasses})$,
  determining the substitution parameters $(\exch^{(\class)}, \eqm^{(\class)})$.
\end{enumerate}
\end{example}

\subsubsection{Link Sequence Concatenation}
\label{sec:concatenation-appendix}

Concatenation decorates a single link with two or more consecutive
sub-links.

\begin{definition}[Link Sequence Concatenation]
\label{def:concat-appendix}
Let $\WCFG$ be a WCFG with nonterminal $X$ representing a link.
The \emph{binary concatenation} of $X$ replaces it with
two new nonterminals $X_A, X_B$ and the production:
\begin{equation}
X \to X_A\;X_B \qquad \text{weight } 1
\end{equation}
$X_A$ and $X_B$ may then be independently elaborated.
For $n$-ary concatenation, $X$ is replaced by
$X \to X_1\;X_2\;\cdots\;X_n$.
\end{definition}

\begin{remark}[Concatenation combined with mixture and fragments]
A powerful pattern combines mixture, fragment, and concatenation
to decorate a link with a variable-length sequence of
categorically typed sub-links:
\begin{enumerate}
\item Apply fragment expansion to the link, creating a geometric
  number of sub-link slots.
\item Apply mixture expansion to each sub-link slot, assigning it
  a categorical type.
\item If correlations between adjacent types are desired,
  replace the i.i.d.\ mixture with an HMM output distribution:
  the type sequence is generated by a hidden Markov model
  whose transition matrix captures adjacency preferences.
\end{enumerate}
This yields a variable number of concatenated sub-links
of varying types with Markovian correlations.
In the TKF Genome, this pattern is used for genomic regions
(coding, noncoding, structural) within a top-level link sequence.
\end{remark}

\begin{example}[TKF Genome: region concatenation]
The TKF Genome's top-level grammar has:
\begin{align*}
\texttt{GENOME} &\to \texttt{REGION}\;\texttt{GENOME} && 1-\kappa_R \\
\texttt{GENOME} &\to \epsilon && \kappa_R \\
\texttt{REGION} &\to \texttt{INTER} && p_N \\
                &\;\;|\;\; \texttt{FWDCDS} && p_G/2 \\
                &\;\;|\;\; \texttt{REVCDS} && p_G/2 \\
                &\;\;|\;\; \texttt{STRUCT} && p_S \\
                &\;\;|\;\; \texttt{CONS} && p_C
\end{align*}
Here $\texttt{REGION}$ is both a mixture expansion (over region types)
and a concatenation point (each region type expands into its own sub-grammar).
\end{example}

\subsubsection{Non-Recursive Nesting}
\label{sec:nonrecursive-appendix}

Non-recursive nesting splices a complete sub-grammar into the
transitions of a mortal link, without introducing bifurcation
or self-reference.  This is how nesting works in MixDom.

\begin{definition}[Non-Recursive Nesting]
\label{def:nonrec-nest-appendix}
Let $\WCFG_{\mathrm{outer}}$ be a link grammar with nonterminal
$\texttt{MOR}$ (mortal link), and let $\WCFG_{\mathrm{inner}}$ be
an independent link grammar with its own start symbol
$\texttt{IMM}_{\mathrm{inner}}$ and parameters.
The \emph{non-recursive nesting} of $\WCFG_{\mathrm{inner}}$
into $\texttt{MOR}$ of $\WCFG_{\mathrm{outer}}$ replaces
each terminal-emitting production of $\texttt{MOR}$ with:
\begin{equation}
\texttt{MOR} \to \texttt{IMM}_{\mathrm{inner}}\;\alpha
\qquad \text{weight } \WF_{\mathrm{orig}}
\end{equation}
where $\alpha$ is the original continuation.
The inner grammar generates a complete sequence of inner links,
each with its own parameters, for every outer mortal link.

More precisely, wherever $\texttt{MOR}$ previously emitted a terminal
symbol $c$, it now expands into the entire inner grammar
$\WCFG_{\mathrm{inner}}$, which itself may generate zero or more
characters.  The inner grammar's $\epsilon$-productions
(zero-length inner sequences) give rise to null states
in the combined grammar.
\end{definition}

\begin{remark}
The key difference from recursive nesting (Section~\ref{sec:recursive-appendix})
is that $\WCFG_{\mathrm{inner}}$ does \emph{not} reference any
nonterminals of $\WCFG_{\mathrm{outer}}$.  There is no possibility
of re-entering the outer grammar from within the inner grammar.
This ensures that the combined grammar generates strings from
a regular language (at each level), rather than a context-free language.
\end{remark}

\begin{example}[MixDom as non-recursive nesting]
The MixDom model nests a Markovian fragment process (inner grammar)
into each mortal link of a TKF91 process (outer grammar).
The outer link grammar has parameters $(\insrate_0, \delrate_0)$ and the inner
grammar, for domain type $\dom$, has parameters
$(\insrate_\dom, \delrate_\dom, \ext^{(\dom)}_{\srcfrag\destfrag}, \classdist_{\dom\frag\class}, \exch^{(\class)}, \eqm^{(\class)})$.
Each outer mortal link, instead of emitting a single character,
expands into a domain sequence governed by the Markovian fragment process.

Since the inner grammar can generate the empty string
(the inner link sequence may have length zero), this nesting
creates null states in the combined Pair HMM.  These must be
eliminated by the procedures of Section~\ref{sec:null-management-appendix}.
\end{example}

\subsubsection{Recursive Nesting}
\label{sec:recursive-appendix}

Recursive nesting, used in the TKF Structure Tree, allows
transitions into a mortal link to spawn a bifurcation:
a new nonterminal whose sub-grammar may reference the
original grammar's nonterminals.

\begin{definition}[Recursive Nesting]
\label{def:rec-nest-appendix}
Let $\WCFG$ be a link grammar with nonterminals including
$\texttt{MOR}$ (mortal link).  Let $\texttt{BIF}$ be a new
nonterminal with its own sub-grammar $\WCFG_{\mathrm{bif}}$
that may reference nonterminals of $\WCFG$ (including
$\texttt{IMM}$ and $\texttt{MOR}$).

The \emph{right recursive nesting} of $\texttt{BIF}$ at $\texttt{MOR}$
replaces each terminal-emitting production of $\texttt{MOR}$ with a
mixture of terminal emission and bifurcation:

\noindent\textbf{Before}:
\begin{align*}
\texttt{MOR} &\to c\;\alpha && \WF_{\mathrm{orig}} \cdot \eqm_c
\end{align*}

\noindent\textbf{After}:
\begin{align*}
\texttt{MOR} &\to c\;\alpha && \WF_{\mathrm{orig}} \cdot s \cdot \eqm_c
  && \text{(terminal link, probability } s\text{)} \\
\texttt{MOR} &\to \texttt{BIF}\;\alpha && \WF_{\mathrm{orig}} \cdot (1-s) \cdot \domdist_{\mathrm{bif}}
  && \text{(bifurcation, probability } 1-s\text{)}
\end{align*}
where $s \in (0,1]$ is the probability that a link is a terminal
(character-emitting) link rather than a nesting point, and
$\domdist_{\mathrm{bif}}$ is a distribution over bifurcation types
if there are multiple $\texttt{BIF}$ variants.

For \emph{left recursive nesting}:
\begin{align*}
\texttt{MOR} &\to \alpha\;\texttt{BIF} && \WF_{\mathrm{orig}} \cdot (1-s) \cdot \domdist_{\mathrm{bif}}
\end{align*}

The sub-grammar $\WCFG_{\mathrm{bif}}$ for $\texttt{BIF}$
defines how the bifurcation expands.  Since it may reference
the start symbol of $\WCFG$ (e.g., $\texttt{IMM}$), the
combined grammar is genuinely recursive: link sequences can
contain nested link sequences of arbitrary depth.
\end{definition}

\begin{remark}
The recursive nesting creates null cycles whenever the nested
sub-grammar can generate the empty string.  These must be
handled by the null-state management procedures
(Section~\ref{sec:null-management-appendix}).
In the TKF Structure Tree, the nullability fixed-point iteration
solves for the probability that each nonterminal generates $\epsilon$.
\end{remark}

\begin{example}[TKF Structure Tree: stems and loops]
\label{ex:structure-tree-appendix}
The TKF Structure Tree has two types of link sequences:
\begin{itemize}
\item \emph{Loop sequences} ($L$): left-emitting links generating
  single nucleotides. Nonterminal rule:
  $L \to c_L\;L$ with CTMC expansion (left emission).
\item \emph{Stem sequences} ($S$): LR-emitting links generating
  basepairs. Nonterminal rule:
  $S \to c_L\;S\;c_R$ with CTMC expansion (LR emission).
\end{itemize}
The recursive nesting works as follows.
Within a loop sequence, a mortal link may either emit a character
(probability $s_L$) or spawn a stem (probability $1-s_L$):
\[
L \to c\;L \quad |\quad S\;L
\]
At the base of a stem (when the self-loop terminates),
the grammar transitions to a loop:
\[
S \to c_L\;S\;c_R \quad |\quad L
\]
This creates the alternating stem--loop structure
characteristic of RNA secondary structure.
The recursion arises because a loop can spawn a stem,
which eventually returns to a loop, which can spawn another stem,
\emph{ad infinitum}.
\end{example}

\begin{example}[TKF92 with Recursive Domains]
The recursive domain model has nonterminals $\texttt{L}_\dom$
(link sequence for domain $\dom$),
$\texttt{A}_\dom$ (aligned component),
$\texttt{I}_\dom$ (inserted component),
$\texttt{D}_\dom$ (deleted component), etc.
The aligned component rule is:
\begin{align*}
\texttt{A}_\dom &\to c && s_\dom \cdot \eqm_{\dom,c}
  \quad\text{(terminal: character)} \\
\texttt{A}_\dom &\to \texttt{L}_{\dom'} && (1-s_\dom) \cdot \domdist_{\dom\dom'}
  \quad\text{(bifurcation: nested link sequence)}
\end{align*}
Since $\texttt{L}_{\dom'}$ can reference $\texttt{L}_\dom$
(if $\dom' = \dom$ or through a chain of domain transitions),
this creates genuine recursion: link sequences contain domains
that contain further link sequences.
\end{example}

\subsubsection{Evolution}
\label{sec:evolution-appendix}

The Evolution elaboration converts a single-sequence grammar
(describing the stationary distribution) into a pair grammar
(describing the joint ancestor--descendant distribution
at evolutionary time $\evoltime$).

\begin{definition}[Evolution Elaboration]
\label{def:evolution-appendix}
Let $\WCFG_1$ be a single-sequence grammar with link nonterminal
$X$ generating terminals from alphabet $\alphabet$ with equilibrium
distribution $\eqm$.  Let $\exch$ be the CTMC rate matrix and
$\evoltime > 0$ the evolutionary time.

Define the TKF parameters:
\begin{align*}
\alpha &= e^{-\delrate \evoltime}, \quad
\beta = \frac{\insrate(e^{-\insrate\evoltime} - e^{-\delrate\evoltime})}
  {\delrate e^{-\insrate\evoltime} - \insrate e^{-\delrate\evoltime}}, \quad
\gamma = 1 - \frac{\delrate\beta}{\insrate(1-\alpha)}, \quad
\kappa = \frac{\insrate}{\delrate}
\end{align*}

The evolution elaboration replaces each link nonterminal $X$
in $\WCFG_1$ with three nonterminals $X_\mat$, $X_\ins$, $X_\del$
in the pair grammar $\WCFG_2$:

\medskip
\noindent\textbf{Before} (single-sequence grammar, left-emitting):
\begin{align*}
X &\to c\;X\;X && \kappa \cdot \eqm_c  && \text{(link with offspring)} \\
X &\to c        && (1-\kappa) \cdot \eqm_c  && \text{(terminal link)}
\end{align*}

\noindent\textbf{After} (pair grammar):

For the \emph{immortal link} continuation, each nonterminal
$Y$ that previously generated the link sequence
$Y \to X\;Y\;|\;\epsilon$ transforms as follows:
\begin{align*}
Y_\mat &\to X_\mat\;Y_\mat &&
  (1-\beta_Y)\kappa\alpha \\
Y_\mat &\to X_\ins\;Y_\mat &&
  \beta_Y \\
Y_\mat &\to X_\del\;Y_\del &&
  (1-\beta_Y)\kappa(1-\alpha) \\
Y_\mat &\to \epsilon &&
  (1-\beta_Y)(1-\kappa) \\[6pt]
Y_\del &\to X_\mat\;Y_\mat &&
  (1-\gamma_Y)\kappa\alpha \\
Y_\del &\to X_\ins\;Y_\mat &&
  \gamma_Y \\
Y_\del &\to X_\del\;Y_\del &&
  (1-\gamma_Y)\kappa(1-\alpha) \\
Y_\del &\to \epsilon &&
  (1-\gamma_Y)(1-\kappa) \\
\end{align*}
where $\beta_Y$ and $\gamma_Y$ use the appropriate $(\insrate, \delrate)$
for the link sequence that $Y$ belongs to.

The subscript indicates the alignment type:
\begin{itemize}
\item $X_\mat$ (\textbf{match}): ancestral link survived; emits
  aligned pair $(c_a, c_d)$ with probability
  $\eqm_{c_a} \cdot \exp(\revsub \evoltime)_{c_a c_d}$.
\item $X_\ins$ (\textbf{insert}): new link born in descendant;
  emits descendant-only character $c_d$ with probability $\eqm_{c_d}$.
\item $X_\del$ (\textbf{delete}): ancestral link died; emits
  ancestor-only character $c_a$ with probability $\eqm_{c_a}$.
\end{itemize}

The transition weights are precisely the entries of the TKF91
Pair HMM transition matrix $\tkftrans(\insrate, \delrate, \evoltime)$.
\end{definition}

\begin{proposition}
The evolution elaboration roughly triples the number of nonterminals
(each single-sequence nonterminal becomes three pair nonterminals).
For LR-emitting nonterminals (as in stem sequences),
match states emit paired tuples $(c_L^a, c_R^a, c_L^d, c_R^d)$,
insert states emit $(c_L^d, c_R^d)$, and delete states emit
$(c_L^a, c_R^a)$.
\end{proposition}

\begin{remark}
When fragment expansion has been applied before evolution,
the fragment self-loop interleaves with the alignment states.
In TKF92, the Pair HMM transition matrix becomes
$\tkftrans'$ where the $\mat$, $\ins$, and $\del$ self-loops
gain a fragment-extension component:
\[
\tkftrans'_{aa} = \ext + (1-\ext)\tkftrans_{aa}
\qquad \text{for } a \in \{\mat, \ins, \del\}
\]
and off-diagonal transitions are scaled by $(1-\ext)$:
\[
\tkftrans'_{ab} = (1-\ext)\tkftrans_{ab}
\qquad \text{for } a \neq b
\]
\end{remark}

\begin{remark}
For LR-emitting grammars (such as stem sequences in the TKF Structure Tree),
the evolution elaboration creates nonterminals $X_M, X_I, X_D$
whose left and right emissions are correlated:
\begin{itemize}
\item $X_M \to c_L^x\;c_L^y\;X_M\;c_R^y\;c_R^x$:
  ancestral basepair $(c_L^x, c_R^x)$ evolved to
  descendant basepair $(c_L^y, c_R^y)$ (match).
\item $X_I$: emits only descendant basepair $c_L^y\;\cdots\;c_R^y$ (insertion).
\item $X_D$: emits only ancestral basepair $c_L^x\;\cdots\;c_R^x$ (deletion).
\end{itemize}
The ancestor terminals go on the outside, the descendant terminals
on the inside (or vice versa), preserving the palindromic nesting.
\end{remark}

\begin{example}[TKF91 Pair HMM]
Applying evolution to the TKF91 single-sequence grammar
(Example~\ref{ex:ctmc-tkf91-appendix}) yields the standard 5-state
Pair HMM $(\sta, \mat, \ins, \del, \fin)$ with transition matrix
$\tkftrans(\insrate, \delrate, \evoltime)$.
\end{example}

\subsection{Null State Management}
\label{sec:null-management-appendix}

\subsubsection{Null State Identification}

\begin{definition}[Null state]
A nonterminal $X$ in a WCFG is \emph{nullable} if there exists
a derivation $X \Rightarrow^* \epsilon$.
The \emph{nullability} $\nullability(X) = P(X \Rightarrow^* \epsilon)$
is the probability that a parse tree rooted at $X$ yields the
empty string.
\end{definition}

Elaborations that create null states include:
\begin{enumerate}
\item \textbf{Non-recursive nesting}: the inner grammar may generate
  the empty string (e.g., a domain in MixDom may contain zero fragments).
\item \textbf{Recursive nesting}: nested link sequences can be empty.
\item \textbf{Mixture expansion} combined with nesting:
  a mixture component that expands into a nullable sub-grammar.
\end{enumerate}

In the MixDom Pair HMM, null states arise because each domain's
inner TKF92 process can generate an empty sequence.
The probability of an empty domain at evolutionary time $\evoltime$
is $\emptyseg_\evoltime = \sum_\dom \domdist_\dom (1-\kappa_\dom)(1-\beta_\dom)$.

\subsubsection{Null State Removal: The $(I - T_{NN})^{-1}$ Closure}
\label{sec:null-removal-appendix}

\begin{definition}[Null Closure]
Let $\WCFG$ be a WCFG (or equivalently an HMM/transducer) with
states partitioned into emitting states $\Omega$ and non-emitting
(null) states $\mathcal{Z}$.  Let $T_{\mathcal{Z}\mathcal{Z}}$
be the submatrix of transition weights among null states.
The \emph{null closure} is:
\begin{equation}
\nullcl = (I - T_{\mathcal{Z}\mathcal{Z}})^{-1}
= \sum_{k=0}^{\infty} T_{\mathcal{Z}\mathcal{Z}}^k
\end{equation}
This converges provided the spectral radius
$\rho(T_{\mathcal{Z}\mathcal{Z}}) < 1$.
\end{definition}

\begin{proposition}[Effective Transition Matrix]
The effective transition matrix between emitting (and start/end) states,
with all null-state paths summed out, is:
\begin{equation}
\effT = T_{\Omega\Omega}
+ T_{\Omega\mathcal{Z}} \cdot (I - T_{\mathcal{Z}\mathcal{Z}})^{-1}
\cdot T_{\mathcal{Z}\Omega}
\end{equation}
where subscripts denote submatrices restricted to the indicated
state sets.
\end{proposition}


\subsubsection{Null Cycle Detection and Removal}
\label{sec:null-cycles-appendix}

\begin{definition}[Null Cycle]
A \emph{null cycle} in a WCFG is a chain of unit productions
(productions whose right-hand side is a single nonterminal)
that returns to the starting nonterminal:
$X \to Y_1 \to Y_2 \to \cdots \to X$.
In the context of HMMs/transducers, this corresponds to a
cycle among non-emitting states.
\end{definition}

Null cycles arise in two main situations:
\begin{enumerate}
\item \textbf{Non-recursive nesting with empty inner grammars}:
  when the inner grammar can produce $\epsilon$, a path
  $\sta \to \texttt{MOR}_{\mathrm{outer}} \to
  \texttt{IMM}_{\mathrm{inner}} \to \epsilon \to
  \texttt{MOR}_{\mathrm{outer}}$ creates a cycle
  through null states.
\item \textbf{Recursive nesting}: the chain
  $\texttt{L}'_\dom \to \texttt{M}'_\dom \to \texttt{S}'_\dom
  \to \texttt{A}'_\dom \to \texttt{L}'_{\dom'}$ creates a
  null cycle when $\dom' = \dom$ or when the chain of domain
  transitions eventually returns to $\dom$.
\end{enumerate}

\begin{definition}[Null Cycle Removal]
To remove null cycles from a WCFG:
\begin{enumerate}
\item \textbf{Compute nullabilities}: for each nonterminal $X$,
  compute $\nullability(X) = P(X \Rightarrow^* \epsilon)$.
  In the non-recursive case, this can be done in closed form.
  In the recursive case, iterate the fixed-point equations:
  \begin{equation}
  \nullability(X) = \sum_{(X \to \alpha) \in \PR}
    \WF(X \to \alpha) \cdot \prod_{Y \in \alpha} \nullability(Y)
  \end{equation}
  where the product is over nonterminals in $\alpha$,
  with $\nullability(\text{terminal}) = 0$ and
  $\nullability(\epsilon) = 1$.

  Initialize $\nullability^{(0)}(X) = 0$ for all $X$
  and iterate until convergence.

\item \textbf{Create non-nullable copies}: for each nullable
  nonterminal $X$, create $X'$ whose productions never generate
  $\epsilon$.  For a bifurcation rule $X \to Y\;Z$,
  the non-nullable version adds:
  \begin{align*}
  X' &\to Y'\;Z' && \WF(X \to YZ) \\
  X' &\to Y' && \WF(X \to YZ) \cdot \nullability(Z) \\
  X' &\to Z' && \WF(X \to YZ) \cdot \nullability(Y)
  \end{align*}
  This accounts for the two ways one child can be null.

\item \textbf{Remove unit-production cycles}: identify cycles
  $X' \to Y'_1 \to \cdots \to X'$ among the non-nullable
  nonterminals.  For each such cycle, compute the transition
  matrix $\mathcal{A}$ among the cycle's nonterminals and replace
  the cycle with its closure $(I - \mathcal{A})^{-1}$,
  distributing the accumulated weight to the non-cyclic
  continuations.
\end{enumerate}
\end{definition}

\begin{example}[Recursive domains: nullability fixed point]
In the recursive domain model,
$\nullability(\texttt{C}_\dom)$ (nullability of the child
link sequence nonterminal) satisfies:
\[
\nullability(\texttt{C}_\dom) =
\frac{1 - \kappa_\dom}
{1 - \kappa_\dom (1-s_\dom) \sum_{\dom'}
  \domdist_{\dom\dom'} \nullability(\texttt{C}_{\dom'})}
\]
This is solved by initializing $x_\dom^{(0)} = 0$ and iterating:
\[
x_\dom^{(k+1)} =
\frac{1 - \kappa_\dom}
{1 - \kappa_\dom (1-s_\dom) \sum_{\dom'}
  \domdist_{\dom\dom'} x_{\dom'}^{(k)}}
\]
After convergence, the null cycles
$\texttt{L}'_\dom \to \texttt{L}'_{\dom'}$ and
$\texttt{C}'_\dom \to \texttt{C}'_{\dom'}$
are removed using the $(I - \mathcal{A})^{-1}$ and
$(I - \mathcal{B})^{-1}$ closures respectively,
where $\mathcal{A}$ and $\mathcal{B}$ are the
$\ndom \times \ndom$ transition matrices among domain types.
\end{example}

\subsection{Composition Properties}

\subsubsection{Commutativity and Order}

The elaboration rules do not, in general, commute.
The following table summarizes the ordering constraints.

\begin{center}
\begin{tabular}{p{4cm}p{8cm}}
\toprule
\textbf{Constraint} & \textbf{Reason} \\
\midrule
CTMC expansion commutes with mixture expansion &
Mixture selects which CTMC parameters to use;
the order of these two operations does not affect the final grammar.
Both orderings produce the same set of productions. \\[6pt]
Fragment expansion must follow (or be simultaneous with) CTMC expansion &
Fragment expansion wraps terminal emissions in a geometric self-loop.
Without terminals, fragment expansion produces fragments of $\epsilon$'s. \\[6pt]
Non-recursive nesting must follow both CTMC and fragment expansion of the inner grammar &
The inner grammar must be fully specified before it can be spliced
into the outer grammar. \\[6pt]
Evolution must be applied last (after all structural elaborations) &
Evolution triples the nonterminals and introduces
alignment-dependent transition weights ($\alpha, \beta, \gamma$).
Applying structural elaborations after evolution would require
modifying all three copies independently. \\[6pt]
Mixture expansion commutes with concatenation &
Both are structural operations on different aspects of a link. \\[6pt]
Null state removal must follow all nullable elaborations but precede distillation &
All null states must be identified before they can be summed out.
The distilled order-1 machines assume null-free grammars. \\
\bottomrule
\end{tabular}
\end{center}

\subsubsection{Validity Conditions}

\begin{definition}[Well-Formed Elaborated Grammar]
An elaborated grammar $\WCFG'$ is \emph{well-formed} if:
\begin{enumerate}
\item \textbf{Properness}: for every nonterminal $X$,
  the production weights sum to 1.
\item \textbf{No unresolved null cycles}: after null state removal,
  no cycles among non-emitting states remain.
  Equivalently, $\rho(T_{\mathcal{Z}\mathcal{Z}}) < 1$
  for the null-state transition matrix.
\item \textbf{Convergent nullability}: for recursive grammars,
  the fixed-point iteration for nullabilities converges.
  A sufficient condition is that $\kappa_\dom < 1$ and
  $s_\dom > 0$ for all domain types (every link has a positive
  probability of being terminal rather than a nesting point).
\item \textbf{Finite expected derivation length}: the expected
  total number of terminals generated is finite.
  For the link grammar, this requires $\kappa < 1$.
  For recursive nesting, additional conditions on the nesting
  probabilities are needed.
\end{enumerate}
\end{definition}

\begin{proposition}
Each elaboration rule defined above is validity-preserving
under its stated conditions.  Composition of validity-preserving
elaborations is validity-preserving, provided the ordering
constraints above are respected.
\end{proposition}

\subsubsection{Derivation of Existing Models}
\label{sec:derivations-appendix}

We now show explicitly how each known TKF-family model arises
as a sequence of elaborations applied to the base link grammar.

\paragraph{TKF91.}

\begin{equation*}
\boxed{\text{TKF91}} =
\WCFG_{\mathrm{link}}(\kappa)
\xrightarrow{\text{CTMC}(\alphabet, \eqm, \exch)}
\WCFG_{\mathrm{TKF91}}
\end{equation*}

\noindent Steps:
\begin{enumerate}
\item Start with the link grammar $\WCFG_{\mathrm{link}}(\kappa)$
  (Definition~\ref{def:link-grammar-appendix}).
\item Apply left CTMC expansion (Definition~\ref{def:ctmc-appendix}) to
  $\texttt{MOR}$ with alphabet $\alphabet$, equilibrium $\eqm$,
  rate matrix $\exch$.
\end{enumerate}
Result: each mortal link emits a single character from $\eqm$.
The pair grammar is obtained by applying evolution
(Definition~\ref{def:evolution-appendix}), yielding the standard 5-state
Pair HMM with transition matrix $\tkftrans(\insrate, \delrate, \evoltime)$.

\paragraph{TKF92.}

\begin{equation*}
\boxed{\text{TKF92}} =
\WCFG_{\mathrm{link}}(\kappa)
\xrightarrow{\text{Frag}(\ext)}
\xrightarrow{\text{CTMC}(\alphabet, \eqm, \exch)}
\WCFG_{\mathrm{TKF92}}
\end{equation*}

\noindent Steps:
\begin{enumerate}
\item Start with $\WCFG_{\mathrm{link}}(\kappa)$.
\item Apply fragment expansion (Definition~\ref{def:fragment-appendix})
  with extension probability $\ext$.
\item Apply left CTMC expansion to each fragment position.
\end{enumerate}
Result: each mortal link emits a fragment of $K \sim \geomdist(\ext)$
characters.  The pair grammar has self-looping match/insert/delete
states with fragment-extension probability $\ext$.

\paragraph{MixDom: Markovian fragments.}

\begin{equation*}
\boxed{\text{MixDom}} =
\WCFG_{\mathrm{link}}(\kappa_0)
\xrightarrow{\text{Mix}(\domdist_\dom)}
\xrightarrow[\text{for each } \dom]{\text{NonRecNest}\bigl(
  \WCFG_{\mathrm{link}}(\kappa_\dom)
  \xrightarrow{\text{HMM}(\ext^{(\dom)})}
  \xrightarrow{\text{Mix}(\fragdist_{\dom\frag})}
  \xrightarrow{\text{Mix}(\classdist_{\dom\frag\class})}
  \xrightarrow{\text{CTMC}(\alphabet, \eqm^{(\class)}, \exch^{(\class)})}
\bigr)}
\WCFG_{\mathrm{MixDom}}
\end{equation*}

\noindent Steps:
\begin{enumerate}
\item Start with $\WCFG_{\mathrm{link}}(\kappa_0)$ (outer/top-level link grammar).
\item Apply mixture expansion (Definition~\ref{def:mixture-appendix}) to
  $\texttt{MOR}$ with domain types $\dom \sim \catdist(\domdist_1,\ldots,\domdist_\ndom)$.
\item For each domain type $\dom$, construct an inner grammar:
  \begin{enumerate}
  \item Start with $\WCFG_{\mathrm{link}}(\kappa_\dom)$ (inner link grammar).
  \item Replace geometric fragment extension with the Markovian fragment HMM
    governed by the $\nfrag \times \nfrag$ transition matrix $\ext^{(\dom)}_{\srcfrag\destfrag}$.
  \item Apply mixture expansion for initial fragment types $\frag \sim \catdist(\fragdist_{\dom 1},\ldots,\fragdist_{\dom\nfrag})$.
  \item Apply mixture expansion for site classes $\class \sim \catdist(\classdist_{\dom\frag 1},\ldots,\classdist_{\dom\frag\nclasses})$.
  \item Apply CTMC expansion with $(\alphabet, \eqm^{(\class)}, \exch^{(\class)})$.
  \end{enumerate}
\item Apply non-recursive nesting (Definition~\ref{def:nonrec-nest-appendix}):
  splice each domain's inner grammar into the corresponding
  outer mortal link.
\item Apply null state removal (Section~\ref{sec:null-removal-appendix}):
  the inner grammar can generate empty sequences
  (null domains), creating the null states $\mnull, \inull, \dnull$
  in the null-separated Pair HMM.  The $(I - T_{\mathcal{Z}\mathcal{Z}})^{-1}$
  closure reduces the 8-state null-separated Pair HMM to the effective
  5-state matrix $\nonemptytrans$.
\end{enumerate}

\paragraph{TKF Structure Tree.}

\begin{equation*}
\boxed{\text{TKF Structure Tree}} =
\WCFG_{\mathrm{link}}(\kappa_L)
\xrightarrow{\text{CTMC}_L(\alphabet, \eqm_L, \exch_L)}
\xrightarrow{\text{RecNest}\bigl(
  \WCFG_{\mathrm{link}}(\kappa_S)
  \xrightarrow{\text{CTMC}_{LR}(\alphabet^2, \eqm_S, \exch_S)}
\bigr)}
\WCFG_{\mathrm{ST}}
\end{equation*}

\noindent Steps:
\begin{enumerate}
\item Start with $\WCFG_{\mathrm{link}}(\kappa_L)$ (loop link grammar).
\item Apply left CTMC expansion for loops
  with $(\alphabet, \eqm_L, \exch_L)$.
\item Apply recursive nesting (Definition~\ref{def:rec-nest-appendix}):
  within the loop, a mortal link may spawn a stem.
  The stem sub-grammar is:
  \begin{enumerate}
  \item Start with $\WCFG_{\mathrm{link}}(\kappa_S)$ (stem link grammar).
  \item Apply LR CTMC expansion with basepair alphabet
    $(\alphabet \times \alphabet, \eqm_S, \exch_S)$.
  \item At stem termination, return to the loop grammar
    (creating the recursion $S \to L$).
  \end{enumerate}
\item The loop grammar now has two types of mortal links:
  \begin{itemize}
  \item Character-emitting links (probability $s_L$):
    emit a single nucleotide $c$ with $\eqm_L(c)$.
  \item Stem-spawning links (probability $1 - s_L$):
    expand into a stem $S$ nonterminal.
  \end{itemize}
\end{enumerate}

The LR emission in the stem grammar is critical:
the rule $S \to c_L\;S\;c_R$ emits characters on
\emph{both} sides of the recursive expansion,
generating the palindromic base-pairing structure of RNA stems.

\paragraph{TKF Genome.}

\begin{equation*}
\boxed{\text{TKF Genome}} =
\WCFG_{\mathrm{link}}(\kappa_R)
\xrightarrow{\text{Concat+Mix(region types)}}
\xrightarrow{\text{various nested elaborations per region type}}
\WCFG_{\mathrm{Genome}}
\end{equation*}

\noindent Steps:
\begin{enumerate}
\item Start with $\WCFG_{\mathrm{link}}(\kappa_R)$ (top-level genomic region grammar).
\item Apply concatenation + mixture: each link is a ``region''
  selected from types
  $\{\texttt{INTER}, \texttt{FWDCDS}, \texttt{REVCDS},
    \texttt{STRUCT}, \texttt{CONS}\}$
  with probabilities $(p_N, p_G/2, p_G/2, p_S, p_C)$.
\item Each region type undergoes its own elaboration chain:
  \begin{itemize}
  \item \texttt{INTER}: link grammar + left CTMC expansion
    (single nucleotides, neutral evolution).
  \item \texttt{FWDCDS}/\texttt{REVCDS}: link grammar + concatenation
    into codons ($c_1 c_2 c_3$) + CTMC expansion
    with codon substitution model + recursive nesting for introns
    (introns contain a nested $\texttt{GENOME}$ nonterminal,
    flanked by splice donor/acceptor sites).
  \item \texttt{STRUCT}: link grammar with LR CTMC expansion (stems)
    + recursive nesting into loops (left CTMC expansion).
  \item \texttt{CONS}: link grammar + left CTMC expansion
    (conserved elements).
  \end{itemize}
\item The intron nesting creates recursion:
  $\texttt{GENOME} \to \texttt{CDS} \to \texttt{CODON} \to
  \texttt{INTRON} \to \texttt{GENOME}$.
\end{enumerate}

\subsection{Toward Implementation}

\subsubsection{Grammar Objects with Transformation Methods}

The elaboration rules defined above can be implemented as methods
on a grammar object:

\begin{verbatim}
class WCFG:
    nonterminals: Set[str]
    terminals: Set[str]
    productions: List[Production]
    start: str

    def ctmc_expand(self, nonterminal, alphabet,
                    equilibrium, emission_type='left'):
        """CTMC Expansion"""

    def fragment_expand(self, nonterminal, extension_prob):
        """Fragment Expansion"""

    def mixture_expand(self, nonterminal, components, weights):
        """Mixture Expansion"""

    def concatenate(self, nonterminal, parts):
        """Link Sequence Concatenation"""

    def nest_nonrecursive(self, nonterminal, inner_grammar):
        """Non-Recursive Nesting"""

    def nest_recursive(self, nonterminal, bifurcation_grammar,
                       terminal_prob, side='right'):
        """Recursive Nesting"""

    def evolve(self, time, rates):
        """Evolution"""

    def remove_null_states(self):
        """Null State Removal"""
\end{verbatim}

Each method validates the preconditions (e.g., that fragment
expansion targets a nonterminal with terminal emissions),
performs the transformation, and returns the modified grammar.

\subsubsection{Automatic Derivation of DP Algorithms}

Given an elaborated grammar, the dynamic programming algorithm
(Forward, Backward, Inside, Viterbi) can be derived automatically
by:

\begin{enumerate}
\item \textbf{State space identification}: each nonterminal in the
  elaborated grammar corresponds to a state in the DP.
  Emitting states correspond to observable positions (sequence characters);
  non-emitting states are handled by null closure.

\item \textbf{Transition structure}: the productions define
  the recurrence relations.  Linear (non-branching) productions
  yield HMM-style recurrences; branching (bifurcation) productions
  yield CYK-style recurrences.

\item \textbf{Emission probabilities}: determined by the CTMC
  parameters and the emission type (left, right, LR, match/insert/delete).

\item \textbf{Fill order}: determined by the topological sort of
  nonterminals (after null cycle removal).
  For recursive grammars, the fill order follows the CYK pattern
  (by span length for context-free rules).
\end{enumerate}

\subsubsection{Connection to Existing Frameworks}

The elaboration rules defined here are compatible with existing
grammar and transducer frameworks:

\begin{itemize}
\item \textbf{Transducer composition}: the Evolution elaboration
  produces a transducer (Pair HMM / WFST).  These can be composed
  on phylogenetic trees using the standard composition and
  intersection operations for Mealy machines in waiting-machine
  normal form.

\item \textbf{SCFG parsers}: the elaborated grammars (especially
  those with recursive nesting and LR emission) are SCFGs
  amenable to Inside/Outside parsing algorithms.

\item \textbf{Distillation}: the elaborated pair grammars can be
  distilled to order-1 machines (HMMs and WFSTs) by computing
  adjacency frequencies and normalizing, as described in
  Section~\ref{sec:distill-hmm}.
  This step loses the hierarchical structure but
  produces compact machines suitable for phylogenetic composition.
\end{itemize}
